\documentclass[11pt,a4paper]{article}

\usepackage[margin=2.5cm]{geometry}
\usepackage{amsmath,amssymb,amsthm}
\usepackage{bm}
\usepackage{booktabs}
\usepackage{xcolor}
\usepackage{hyperref}
\usepackage{enumitem}

\newtheorem{remark}{Remark}

\newcommand{\R}{\mathbb{R}}
\newcommand{\dom}{\Omega}
\newcommand{\pore}[1]{\Omega_{#1}}
\newcommand{\iface}{\gamma}
\newcommand{\vel}{\bm{u}}
\newcommand{\pres}{p}
\newcommand{\stress}{\bm{\sigma}}
\newcommand{\normal}{\bm{n}}
\newcommand{\trac}{\bm{\lambda}}
\newcommand{\DtN}{\mathbf{G}}
\newcommand{\Schur}{\mathbf{S}}
\newcommand{\one}{\mathbf{1}}
\newcommand{\zero}{\bm{0}}
\newcommand{\modeset}{\mathcal{M}}

\hypersetup{colorlinks=true,linkcolor=blue,urlcolor=blue,citecolor=blue}

\begin{document}

\title{\textbf{Affine-DDPNM: Methodology}\\[4pt]
       \large Version 3 --- revised 2026-08-06}
\author{}
\date{}
\maketitle

We extend the domain-decomposition pore-network method (DD-PNM) of
\cite{sun2025ddpnm} by enriching the interface traction space from a single
constant-normal mode to a nine-mode affine representation.

% ======================================================================
\section{Model problem \& sign convention}
\label{sec:model}
% ======================================================================

Steady incompressible Stokes flow in $\dom \subset \R^3$:
\begin{align}
  -\nu\Delta\vel + \nabla\pres &= \bm{0},                && \text{in } \dom,\\
  \nabla\cdot\vel &= 0,                                    && \text{in } \dom,\\
  \vel &= \bm{0},                                          && \text{on } \Gamma_{\text{wall}},\\
  -\stress(\vel,\pres)\,\normal &= p_b\,\normal,          && \text{on } \Gamma_{\text{in}}\cup\Gamma_{\text{out}},\\
  (I-\normal\otimes\normal)\,\stress(\vel,\pres)\normal &= \bm{0},
                                                          && \text{on } \Gamma_{\text{in}}\cup\Gamma_{\text{out}},
\end{align}
where $\stress(\vel,\pres) := -\pres\mathbf{I} + \nu\nabla\vel$ is the
pseudostress and $p_b = p_{\text{in}}$ on $\Gamma_{\text{in}}$,
$p_b = p_{\text{out}}$ on $\Gamma_{\text{out}}$.
The inlet/outlet condition \textbf{prescribes the normal pseudotraction}
$-\sigma n = p_b n$, not the pressure directly; the distinction matters
when $\nu\nabla\vel$ contributes to the normal stress.

% ======================================================================
\section{Domain partition}
\label{sec:partition}
% ======================================================================

$\dom$ is partitioned into $N$ non-overlapping subdomains
$\{\pore{i}\}_{i=1}^N$.
Interface $\iface_k = \partial\pore{i} \cap \partial\pore{j}$
($k=1,\dots,m$) is assumed \textbf{planar} throughout the theoretical
development (see Remark~\ref{rem:planar} for watershed).
Each $\pore{i}$ has $m_i$ ports, $m_i^u$ internal and $m_i^k$ boundary.
Boolean restriction maps $R_i^u,R_i^k$ follow~\cite{sun2025ddpnm}.

For the Voronoi partition, each interface is the perpendicular-bisector
plane of the two sphere centres; for unequal radii this differs from the
exact clearance saddle surface.  The geometric error is quantified in the
numerical companion.

% ======================================================================
\section{Reference FE discretization}
\label{sec:fe}
% ======================================================================

Taylor--Hood $[P_2]^3$--$P_1$ on body-fitted tetrahedral meshes.
Velocity space $\mathbf{V}_h$, pressure space $Q_h = P_1$.
On $\pore{i}$ with Neumann data:
\begin{equation}\label{eq:local-fe}
  \begin{bmatrix} A_i & B_i^\top \\ B_i & 0 \end{bmatrix}
  \begin{bmatrix} \mathbf{U}_i \\ \mathbf{P}_i \end{bmatrix}
  = \begin{bmatrix} \sum_{r=1}^{m_i} \bm{f}_i^{(r)} \\ \bm{0} \end{bmatrix}.
\end{equation}
No pressure stabilisation is used.

% ======================================================================
\section{Interface traction Ansatz}
\label{sec:affine-ansatz}
% ======================================================================

\subsection{Classic DD-PNM}
\begin{equation}\label{eq:classic}
  -\stress(\vel,\pres)\normal_\gamma = p_\gamma\,\normal_\gamma,
\end{equation}
one scalar $p_\gamma$ per interface.

\subsection{Affine-DDPNM (nine modes per interface)}

For each planar interface $\gamma$, fix an orthonormal frame
$\{\normal_\gamma,\bm{t}_{\gamma,1},\bm{t}_{\gamma,2}\}$ with
$\normal_\gamma$ pointing from the first to the second pore of the pair.
Half-span--normalised coordinates:
\begin{equation}\label{eq:coords}
  s(\bm{x}) = \frac{(\bm{x}-\bm{c}_\gamma)\cdot\bm{t}_{\gamma,1}}{a_\gamma},\qquad
  t(\bm{x}) = \frac{(\bm{x}-\bm{c}_\gamma)\cdot\bm{t}_{\gamma,2}}{b_\gamma},
\end{equation}
with $a_\gamma,b_\gamma$ the half-spans, guaranteeing $s,t=O(1)$ under
shape-regularity.

Signed direction vectors for pore $i$:
\begin{equation}\label{eq:signed-dir}
  \bm{d}_{i,\gamma,\normal} = \normal_{i,\gamma},\qquad
  \bm{d}_{i,\gamma,\bm{t}_\nu} = \sigma_{i,\gamma}\,\bm{t}_{\gamma,\nu}
  \;\;(\nu=1,2),
\end{equation}
where $\sigma_{i,\gamma}=+1$ ($-1$) if $i$ is the first (second) member.
These satisfy $\bm{d}_{j,\gamma,\beta} = -\bm{d}_{i,\gamma,\beta}$.

The \textbf{applied traction} on $\gamma$ as seen from $\pore{i}$ is
\begin{equation}\label{eq:affine-traction}
  -\stress(\vel,\pres)\,\normal_{i,\gamma}
  = \sum_{(\alpha,\beta)\in\modeset}
    \lambda_{\gamma,\alpha,\beta}\;
    \varphi_\alpha(s,t)\,\bm{d}_{i,\gamma,\beta},
\end{equation}
where $\modeset = \{0,s,t\}\times\{\normal,\bm{t}_1,\bm{t}_2\}$,
$\varphi_0=1$, $\varphi_s=s$, $\varphi_t=t$.
The minus sign matches the classic convention~\eqref{eq:classic}.
Nine coefficients $\trac_\gamma \in \R^9$ per interface.

\begin{remark}[Planar interface requirement]\label{rem:planar}
  The theory assumes planar interfaces with fixed $\normal_\gamma$.
  Watershed partitions produce bent interfaces where the mesh-level
  facet normal varies.  In the code, watershed interfaces use
  per-facet $\normal(x)$ for the normal load and a fixed averaged frame
  for tangent directions; this deviates from the nine-mode space and is
  one factor behind the degraded Affine accuracy on watershed geometries.
\end{remark}

% ======================================================================
\section{Local compliance map}
\label{sec:local-dtn}
% ======================================================================

\subsection{Primitive response loads}

For mode $(\alpha,\beta)$ on internal interface $\gamma$ of $\pore{i}$:
\begin{equation}\label{eq:load}
  [\bm{f}_{i,\gamma}^{(\alpha,\beta)}]_j
  = \int_{\gamma} \varphi_\alpha(s,t)\;
    \bm{d}_{i,\gamma,\beta} \cdot \bm{\phi}_j^{(i)} \;\mathrm{d}S.
\end{equation}
With the sign convention~\eqref{eq:affine-traction}, a positive
$\lambda$ acts into the pore; the load vector carries no leading minus
sign---it is the right-hand side of~\eqref{eq:local-fe} directly.
For boundary ports only $(\alpha,\beta)=(0,\normal)$ is used with a
prescribed coefficient.

Solve for each unit load:
\begin{equation}\label{eq:response}
  \begin{bmatrix} A_i & B_i^\top \\ B_i & 0 \end{bmatrix}
  \begin{bmatrix} \mathbf{U}_{i,\gamma}^{(\alpha,\beta)} \\
                   \mathbf{X}_{i,\gamma}^{(\alpha,\beta)} \end{bmatrix}
  = \begin{bmatrix} \bm{f}_{i,\gamma}^{(\alpha,\beta)} \\ \bm{0} \end{bmatrix}.
\end{equation}

\subsection{Compliance map and weighted velocity moments}

$n_i = 9m_i^u + m_i^k$ local coefficients ($m_i^k$ of them prescribed).
$L_i \in \R^{n_{i,u}\times n_i}$: columns = all $\bm{f}_i^{(\mu)}$.
Velocity: $\mathbf{U}_i = \sum_{\gamma\in\mathcal{U}_i}\sum_{\alpha,\beta}
  \lambda_{i,\gamma,\alpha,\beta}\mathbf{U}_{i,\gamma}^{(\alpha,\beta)}
  + \sum_{\gamma\in\mathcal{K}_i}\lambda_{i,\gamma}^k
    \mathbf{U}_{i,\gamma}^{(0,\normal)}$.

\textbf{Weighted velocity moments} on interface $\gamma$:
\begin{equation}\label{eq:moment}
  M_{i,\gamma}^{(\alpha,\beta)}
  = \int_{\gamma} \varphi_\alpha(s,t)\;
    \vel_{i,h} \cdot \bm{d}_{i,\gamma,\beta} \;\mathrm{d}S.
\end{equation}
$M_{i,\gamma}^{(0,\normal)}$ = volumetric flow rate (the only quantity
with the physical meaning of a flux).  $M^{(s,\normal)},M^{(t,\normal)}$
are linear normal flux moments; the six $M^{(\alpha,\bm{t}_\nu)}$ are
tangential velocity moments.  Collectively: \textbf{nine weighted velocity
moments}.

The compliance (Neumann-to-Dirichlet) map:
\begin{equation}\label{eq:dtn}
  \DtN_i = L_i^\top H_i L_i \in \R^{n_i\times n_i},\qquad
  H_i = A_i^{-1} - A_i^{-1}B_i^\top M_i^{-1}B_i A_i^{-1}.
\end{equation}
$\DtN_i \succeq 0$, symmetric.  Partitioned into $\DtN_i^{uu},\DtN_i^{uk}$.

% ======================================================================
\section{Global Schur complement assembly}
\label{sec:global-affine}
% ======================================================================

Enforcing continuity of \textbf{all nine weighted velocity moments}
($M_{i,\gamma}^{(\alpha,\beta)} + M_{j,\gamma}^{(\alpha,\beta)} = 0$):
\begin{equation}\label{eq:schur}
  \Schur\,\trac = \bm{F},\qquad
  \Schur = \sum_{i=1}^N (R_i^u)^\top \DtN_i^{uu} R_i^u,\qquad
  \bm{F} = -\sum_{i=1}^N (R_i^u)^\top \DtN_i^{uk} R_i^k \trac^k.
\end{equation}
$\Schur \in \R^{9m\times 9m}$, symmetric by construction.

\begin{remark}[Conditions for SPD]\label{rem:spd}
  $\Schur$ is positive definite under the following conditions:
  \begin{enumerate}[nosep,leftmargin=*]
    \item \textbf{Global anchoring:} every connected component of the pore
      adjacency graph contains at least one inlet/outlet port;
    \item \textbf{Local coercivity:} every pore has a positive-measure
      no-slip wall segment;
    \item \textbf{Non-degenerate loads:}
      $\operatorname{rank}(H_i^{1/2}L_i) = n_i - 1$ for every pore
      (verified numerically across all tested geometries).
  \end{enumerate}
  Conditions (1) and (2) are distinct: (2) ensures $K_i$ is invertible;
  (1) eliminates the global constant-pressure kernel mode of $\Schur$.
  In the watershed partition, \emph{floating basins} (pores lacking a
  solid-wall facet) violate condition (2) and are merged into neighbours
  before assembly.  This restores local coercivity but does not supply
  condition (1); the inlet/outlet ports independently anchor the global
  system.
  \textbf{Per-face linear independence} of the nine primitive loads is
  necessary but not sufficient for condition~(3), which additionally
  requires the absence of inter-interface linear dependencies after
  projection through $H_i^{1/2}$.
\end{remark}

% ======================================================================
\section{Solution algorithm}
\label{sec:algorithm}
% ======================================================================

\begin{enumerate}[nosep,leftmargin=*,font=\bfseries]
  \item[\bfseries 1.~Offline library] (per pore, in parallel)
    \begin{enumerate}[nosep,leftmargin=*]
      \item Assemble $A_i,B_i$.
      \item For each internal $\gamma$ and $(\alpha,\beta)\in\modeset$:
        assemble $\bm{f}_{i,\gamma}^{(\alpha,\beta)}$,
        solve for $(\mathbf{U},\mathbf{X})$.
      \item For each boundary $\gamma$: solve $(0,\normal)$ load.
      \item Form $\DtN_i$, accumulate $\Schur$, $\bm{F}$.
    \end{enumerate}
  \item[\bfseries 2.~Global solve] $\Schur\trac = \bm{F}$ (sparse direct).
  \item[\bfseries 3.~Reconstruction] (in parallel)
    \begin{align*}
      \mathbf{U}_i &= \sum_{\gamma\in\mathcal{U}_i}\sum_{\alpha,\beta}
        \lambda_{i,\gamma,\alpha,\beta}\mathbf{U}_{i,\gamma}^{(\alpha,\beta)}
        + \sum_{\gamma\in\mathcal{K}_i}\lambda_{i,\gamma}^k
        \mathbf{U}_{i,\gamma}^{(0,\normal)},\\
      \mathbf{P}_i &= \sum_{\gamma\in\mathcal{U}_i}\sum_{\alpha,\beta}
        \lambda_{i,\gamma,\alpha,\beta}\mathbf{X}_{i,\gamma}^{(\alpha,\beta)}
        + \sum_{\gamma\in\mathcal{K}_i}\lambda_{i,\gamma}^k
        \mathbf{X}_{i,\gamma}^{(0,\normal)}.
    \end{align*}
\end{enumerate}

% ======================================================================
\section{Relation to the classic DD-PNM}
\label{sec:reduction}
% ======================================================================

The classic traction space is a subspace of the affine space
($\mathcal{T}_\gamma^{\text{P0}} \subset \mathcal{T}_\gamma^{\text{Affine}}$).
The classic DD-PNM is recovered by \textbf{restricting both trial and test
spaces}:
\[
  \Schur_{\text{P0}} = \mathcal{P}_{\text{P0}}^\top \Schur\,
                       \mathcal{P}_{\text{P0}},\qquad
  \bm{F}_{\text{P0}} = \mathcal{P}_{\text{P0}}^\top \bm{F},
\]
where $\mathcal{P}_{\text{P0}}$ embeds the $m$ classic coefficients into
$\R^{9m}$ by placing them in the $(0,\normal)$ slots.
The classic solution $\bm{p} = \Schur_{\text{P0}}^{-1}\bm{F}_{\text{P0}}$
generally differs from the $(0,\normal)$-components of
$\trac = \Schur^{-1}\bm{F}$: the eight additional test-space constraints
alter the Galerkin projection.  The enrichment guarantee is that
$\operatorname{Range}(\mathcal{P}_{\text{P0}}) \subset
\operatorname{Range}(\mathcal{P}_{\text{aff}})$, so the affine error in
the energy (semi-)norm is never larger than the classic error
(see the companion mathematical properties document).

% ======================================================================
\section{Computational cost}
\label{sec:cost}
% ======================================================================

$n_i \approx 9m_i$ (vs.\ $m_i$ classic).  Offline: $n_i$ independent
local Stokes solves (embarrassingly parallel).  Online: sparse LU of
$\Schur \in \R^{9m\times 9m}$, with at most
$9 \times (\text{max coordination}) \times 9$ non-zero blocks per row.
For multi-boundary-condition sweeps the offline library is built once.

\begin{thebibliography}{99}
\bibitem{sun2025ddpnm}
S.~Sun, Z.~Wang, L.~Zhang, J.~Zhao.
\newblock A domain decomposition approach to pore-network modeling of porous
  media flow.
\newblock \textit{arXiv:2510.13429v2}, 2026.
\end{thebibliography}

\end{document}
